% META: {"id": "TSPE-SUITLIM-CO-026", "chapitre": "TSPE-SUITES-LIMITES", "type_objet": "corrige", "exercice_ref": "TSPE-SUITLIM-EX-026", "capacites_codes": ["C4"], "sources_inspiration": [], "status": "generated"}
% BEGIN-VERIFY
% from sympy import *
% n = symbols('n', integer=True, nonnegative=True)
% u = n + Rational(1,2)*(-1)**n
% assert float(u.subs(n, 10000)) > 9000
% END-VERIFY
\begin{corrige}{TSPE-SUITLIM-EX-026}

\textbf{Q1.} $u_0=-1/2$, $u_1=3/2$, $u_2=3/2$, $u_3=5/2$. Non monotone ($u_0<u_1$ mais $u_1=u_2$).

\textbf{Q2.} $(-1)^n \geq -1$, donc $u_n = n+(-1)^n/2 \geq n-1/2$.

\textbf{Q3.} $n-1/2 \to +\infty$. Par comparaison, $\boxed{\lim u_n = +\infty}$.

\textbf{Q4.} Non, le theoreme exige la croissance. Ici on utilise directement la comparaison.

\end{corrige}
